% This is template.tex of https://github.com/Willie169/LaTeX-ToolKit, licensed under either GPL-3.0-or-later, CC BY-SA 4.0-or-later, or LPPL-1.3c-or-later.
\IfPackageLoadedTF{geometry}{}{\usepackage[margin=1.27cm,foot=0.77cm]{geometry}}
\ifcsname SetMathFont\endcsname\else\newcommand{\SetMathFont}{\setmathfont{XITS Math}\setmathfont[range={cal,bfcal},StylisticSet=1]{XITS Math}}\fi
\ifcsname c@CJK\endcsname\else\newcounter{CJK}\setcounter{CJK}{1}\fi
\ifcsname c@markdown\endcsname\else\newcounter{markdown}\setcounter{markdown}{0}\fi
\ifcsname c@Fonts\endcsname\else\newcounter{Fonts}\setcounter{Fonts}{1}\fi
\ifcsname c@CJKFonts\endcsname\else\newcounter{CJKFonts}\setcounter{CJKFonts}{1}\fi
\ifcsname c@CJKLanguage\endcsname\else\newcounter{CJKLanguage}\setcounter{CJKLanguage}{0}\fi
\ifcsname c@NotoCJKFamily\endcsname\else\newcounter{NotoCJKFamily}\setcounter{NotoCJKFamily}{0}\fi
\ifcsname c@SectionLanguage\endcsname\else\newcounter{SectionLanguage}\setcounter{SectionLanguage}{0}\fi
\ifcsname c@ZhRenew\endcsname\else\newcounter{ZhRenew}\setcounter{ZhRenew}{0}\fi
\ifcsname PhysicsPatch\endcsname\else\newcommand{\PhysicsPatch}{\usepackage[forcedefaultall]{physics-patch}}\fi
\ifnum\value{CJK}=1
\ifx\luatexversion\undefined
\usepackage{fontspec}
\usepackage[CJKspace]{xeCJK}
\else
\usepackage{luaotfload,luatexja-fontspec}
\fi
\usepackage{zhnumber}
\fi
\usepackage{etoolbox,ifthen,xparse,xpatch,xstring}
\usepackage{needspace,parskip,setspace,titlesec,titling,titletoc}
\usepackage[defaultlines=8]{nowidow}
\usepackage{amsmath,amssymb,amsthm,centernot,dsfont,esint,mathrsfs,mathtools,metalogo,stmaryrd,tipa,textcomp,textgreek,upgreek,wasysym,yfonts}
\usepackage{array,booktabs,caption,cellspace,colortbl,diagbox,enumitem,float,graphicx,hanging,longtable,ltxtable,makecell,mdframed,multirow,nicematrix,placeins,ragged2e,subcaption,tabu,xcolor}
\usepackage{braket,chemfig,csquotes,csvsimple,derivative,listings,listingsutf8,pdfpages,verbatim}
\usepackage{tikz}
\usetikzlibrary{circuits.ee.IEC,circuits.logic.US,circuits.logic.IEC}
\usepackage{pgfplots}
\pgfplotsset{compat=1.18} 
\usepackage{siunitx}
\usepackage[siunitx]{circuitikz}
\usepackage[version=4]{mhchem}
\usepackage{arydshln}
\ifnum\value{markdown}=1
\usepackage{markdown}
\markdownSetup{autoIdentifiers=true,blankBeforeBlockquote=true,blankBeforeCodeFence=true,blankBeforeDivFence=true,blankBeforeHeading=true,blankBeforeList=true, bracketedSpans=true,citations=true,definitionLists=true,fancyLists=true,fencedCodeAttributes=true,fencedDivs=true,gfmAutoIdentifiers=true,hashEnumerators=true,headerAttributes=true,inlineCodeAttributes=true,inlineNotes=true,jekyllData=true,linkAttributes=true,lineBlocks=true,mark=true,notes=true,pipeTables=true,preserveTabs=false,rawAttribute=true,smartEllipses=true,strikeThrough=true,stripIndent=true,subscripts=true,superscripts=true,tableAttributes=true,tableCaptions=true,taskLists=true,texMathDollars=true,texMathSingleBackslash=true}
\else\ifnum\value{markdown}=2
\usepackage{markdown}
\markdownSetup{blankBeforeBlockquote=true,blankBeforeCodeFence=true,blankBeforeDivFence=true,blankBeforeHeading=true,blankBeforeList=true, pipeTables=true,preserveTabs=false,smartEllipses=true,strikeThrough=true,stripIndent=true,subscripts=true,superscripts=true,taskLists=true,texMathDollars=true,texMathSingleBackslash=true}
\else\ifnum\value{markdown}=3
\usepackage{markdown}
\markdownSetup{blankBeforeBlockquote=true,blankBeforeCodeFence=true,blankBeforeDivFence=true,blankBeforeHeading=true,blankBeforeList=true, preserveTabs=false,smartEllipses=true,stripIndent=true}
\else\ifnum\value{markdown}=4
\usepackage{markdown}
\markdownSetup{autoIdentifiers=true,bracketedSpans=true,citations=true,definitionLists=true,fancyLists=true,fencedCodeAttributes=true,fencedDivs=true,gfmAutoIdentifiers=true,hashEnumerators=true,headerAttributes=true,inlineCodeAttributes=true,inlineNotes=true,jekyllData=true,linkAttributes=true,lineBlocks=true,mark=true,notes=true,pipeTables=true,preserveTabs=false,rawAttribute=true,smartEllipses=true,strikeThrough=true,stripIndent=true,subscripts=true,superscripts=true,tableAttributes=true,tableCaptions=true,taskLists=true,texMathDollars=true,texMathSingleBackslash=true}
\else\ifnum\value{markdown}=5
\usepackage{markdown}
\markdownSetup{pipeTables=true,preserveTabs=false,smartEllipses=true,strikeThrough=true,stripIndent=true,subscripts=true,superscripts=true,taskLists=true,texMathDollars=true,texMathSingleBackslash=true}
\else\ifnum\value{markdown}=6
\usepackage{markdown}
\markdownSetup{preserveTabs=false,smartEllipses=true,stripIndent=true}
\ifnum\value{markdown}=7
\usepackage{markdown}
\markdownSetup{autoIdentifiers=true,blankBeforeBlockquote=true,blankBeforeCodeFence=true,blankBeforeDivFence=true,blankBeforeHeading=true,blankBeforeList=true, bracketedSpans=true,citations=true,definitionLists=true,fancyLists=true,fencedCodeAttributes=true,fencedDivs=true,gfmAutoIdentifiers=true,hashEnumerators=true,headerAttributes=true,inlineCodeAttributes=true,inlineNotes=true,jekyllData=true,linkAttributes=true,lineBlocks=true,mark=true,notes=true,pipeTables=true,rawAttribute=true,strikeThrough=true,subscripts=true,superscripts=true,tableAttributes=true,tableCaptions=true,taskLists=true,texMathDollars=true,texMathSingleBackslash=true}
\else\ifnum\value{markdown}=8
\usepackage{markdown}
\markdownSetup{blankBeforeBlockquote=true,blankBeforeCodeFence=true,blankBeforeDivFence=true,blankBeforeHeading=true,blankBeforeList=true, pipeTables=true,subscripts=true,superscripts=true,taskLists=true,texMathDollars=true,texMathSingleBackslash=true}
\else\ifnum\value{markdown}=9
\usepackage{markdown}
\markdownSetup{blankBeforeBlockquote=true,blankBeforeCodeFence=true,blankBeforeDivFence=true,blankBeforeHeading=true,blankBeforeList=true}
\else\ifnum\value{markdown}=10
\usepackage{markdown}
\markdownSetup{autoIdentifiers=true,bracketedSpans=true,citations=true,definitionLists=true,fancyLists=true,fencedCodeAttributes=true,fencedDivs=true,gfmAutoIdentifiers=true,hashEnumerators=true,headerAttributes=true,inlineCodeAttributes=true,inlineNotes=true,jekyllData=true,linkAttributes=true,lineBlocks=true,mark=true,notes=true,pipeTables=true,rawAttribute=true,strikeThrough=true,subscripts=true,superscripts=true,tableAttributes=true,tableCaptions=true,taskLists=true,texMathDollars=true,texMathSingleBackslash=true}
\else\ifnum\value{markdown}=11
\usepackage{markdown}
\markdownSetup{pipeTables=true,strikeThrough=true,subscripts=true,superscripts=true,taskLists=true,texMathDollars=true,texMathSingleBackslash=true}
\else\ifnum\value{markdown}=12
\usepackage{markdown}
\fi\fi\fi\fi\fi\fi\fi\fi\fi\fi\fi\fi
\ifcsname PreambleMiddle\endcsname\PreambleMiddle\fi
\usepackage[hidelinks,bookmarksnumbered=true,pdfstartview=FitH]{hyperref}
\usepackage{xurl}
\usepackage{tcolorbox}
\tcbuselibrary{breakable,documentation,fitting,hooks,listings,listingsutf8,poster,raster,skins,theorems}
\usepackage{unicode-math}
\setmathfont{XITS Math}\setmathfont[range={cal,bfcal},StylisticSet=1]{XITS Math}
\ifnum\value{Fonts}=1
\setmainfont{TeX Gyre Heros}[Ligatures=TeX]
\setsansfont{TeX Gyre Termes}[Ligatures=TeX]
\setmonofont{TeX Gyre Cursor}[Ligatures=TeX]
\else\ifnum\value{Fonts}=2
\setmainfont{TeX Gyre Termes}[Ligatures=TeX]
\setsansfont{TeX Gyre Heros}[Ligatures=TeX]
\setmonofont{TeX Gyre Cursor}[Ligatures=TeX]
\else\ifnum\value{Fonts}=3
\setmainfont{TeX Gyre Heros}[Ligatures=TeX]
\setmonofont{TeX Gyre Cursor}[Ligatures=TeX]
\else\ifnum\value{Fonts}=4
\setmainfont{TeX Gyre Termes}[Ligatures=TeX]
\setmonofont{TeX Gyre Cursor}[Ligatures=TeX]
\else\ifnum\value{Fonts}=5
\setmainfont{TeX Gyre Heros}[Ligatures=TeX]
\setsansfont{TeX Gyre Termes}[Ligatures=TeX]
\else\ifnum\value{Fonts}=6
\setmainfont{TeX Gyre Termes}[Ligatures=TeX]
\setsansfont{TeX Gyre Heros}[Ligatures=TeX]
\else\ifnum\value{Fonts}=7
\setmainfont{TeX Gyre Heros}[Ligatures=TeX]
\else\ifnum\value{Fonts}=8
\setmainfont{TeX Gyre Termes}[Ligatures=TeX]
\fi\fi\fi\fi\fi\fi\fi\fi
\ifnum\value{CJK}=1
\ifx\luatexversion\undefined
\ifnum\value{CJKFonts}=0\else
\ifnum\value{CJKLanguage}=0
\ifnum\value{CJKFonts}=1
\setCJKmainfont{Noto Sans CJK TC}[Ligatures=TeX]
\setCJKsansfont{Noto Serif CJK TC}[Ligatures=TeX]
\setCJKmonofont{Noto Sans CJK MonoCJKtc}[Ligatures=TeX]
\else\ifnum\value{CJKFonts}=2
\setCJKmainfont{Noto Serif CJK TC}[Ligatures=TeX]
\setCJKsansfont{Noto Sans CJK TC}[Ligatures=TeX]
\setCJKmonofont{Noto Sans CJK MonoCJKtc}[Ligatures=TeX]
\else\ifnum\value{CJKFonts}=3
\setCJKmainfont{Noto Sans CJK TC}[Ligatures=TeX]
\setCJKmonofont{Noto Sans CJK MonoCJKtc}[Ligatures=TeX]
\else\ifnum\value{CJKFonts}=4
\setCJKmainfont{Noto Serif CJK TC}[Ligatures=TeX]
\setCJKmonofont{Noto Sans CJK MonoCJKtc}[Ligatures=TeX]
\else\ifnum\value{CJKFonts}=5
\setCJKmainfont{Noto Sans CJK TC}[Ligatures=TeX]
\setCJKsansfont{Noto Serif CJK TC}[Ligatures=TeX]
\else\ifnum\value{CJKFonts}=6
\setCJKmainfont{Noto Serif CJK TC}[Ligatures=TeX]
\setCJKsansfont{Noto Sans CJK TC}[Ligatures=TeX]
\else\ifnum\value{CJKFonts}=7
\setCJKmainfont{Noto Sans CJK TC}[Ligatures=TeX]
\else\ifnum\value{CJKFonts}=8
\setCJKmainfont{Noto Serif CJK TC}[Ligatures=TeX]
\fi\fi\fi\fi\fi\fi\fi\fi\else
\ifnum\value{CJKLanguage}=1
\ifnum\value{CJKFonts}=1
\setCJKmainfont{Noto Sans CJK SC}[Ligatures=TeX]
\setCJKsansfont{Noto Serif CJK SC}[Ligatures=TeX]
\setCJKmonofont{Noto Sans CJK MonoCJKsc}[Ligatures=TeX]
\else\ifnum\value{CJKFonts}=2
\setCJKmainfont{Noto Serif CJK SC}[Ligatures=TeX]
\setCJKsansfont{Noto Sans CJK SC}[Ligatures=TeX]
\setCJKmonofont{Noto Sans CJK MonoCJKsc}[Ligatures=TeX]
\else\ifnum\value{CJKFonts}=3
\setCJKmainfont{Noto Sans CJK SC}[Ligatures=TeX]
\setCJKmonofont{Noto Sans CJK MonoCJKsc}[Ligatures=TeX]
\else\ifnum\value{CJKFonts}=4
\setCJKmainfont{Noto Serif CJK SC}[Ligatures=TeX]
\setCJKmonofont{Noto Sans CJK MonoCJKsc}[Ligatures=TeX]
\else\ifnum\value{CJKFonts}=5
\setCJKmainfont{Noto Sans CJK SC}[Ligatures=TeX]
\setCJKsansfont{Noto Serif CJK SC}[Ligatures=TeX]
\else\ifnum\value{CJKFonts}=6
\setCJKmainfont{Noto Serif CJK SC}[Ligatures=TeX]
\setCJKsansfont{Noto Sans CJK SC}[Ligatures=TeX]
\else\ifnum\value{CJKFonts}=7
\setCJKmainfont{Noto Sans CJK SC}[Ligatures=TeX]
\else\ifnum\value{CJKFonts}=8
\setCJKmainfont{Noto Serif CJK SC}[Ligatures=TeX]
\fi\fi\fi\fi\fi\fi\fi\fi\else
\ifnum\value{CJKLanguage}=2
\ifnum\value{CJKFonts}=1
\setCJKmainfont{Noto Sans CJK HK}[Ligatures=TeX]
\setCJKsansfont{Noto Serif CJK HK}[Ligatures=TeX]
\setCJKmonofont{Noto Sans CJK MonoCJKhk}[Ligatures=TeX]
\else\ifnum\value{CJKFonts}=2
\setCJKmainfont{Noto Serif CJK HK}[Ligatures=TeX]
\setCJKsansfont{Noto Sans CJK HK}[Ligatures=TeX]
\setCJKmonofont{Noto Sans CJK MonoCJKhk}[Ligatures=TeX]
\else\ifnum\value{CJKFonts}=3
\setCJKmainfont{Noto Sans CJK HK}[Ligatures=TeX]
\setCJKmonofont{Noto Sans CJK MonoCJKhk}[Ligatures=TeX]
\else\ifnum\value{CJKFonts}=4
\setCJKmainfont{Noto Serif CJK HK}[Ligatures=TeX]
\setCJKmonofont{Noto Sans CJK MonoCJKhk}[Ligatures=TeX]
\else\ifnum\value{CJKFonts}=5
\setCJKmainfont{Noto Sans CJK HK}[Ligatures=TeX]
\setCJKsansfont{Noto Serif CJK HK}[Ligatures=TeX]
\else\ifnum\value{CJKFonts}=6
\setCJKmainfont{Noto Serif CJK HK}[Ligatures=TeX]
\setCJKsansfont{Noto Sans CJK HK}[Ligatures=TeX]
\else\ifnum\value{CJKFonts}=7
\setCJKmainfont{Noto Sans CJK HK}[Ligatures=TeX]
\else\ifnum\value{CJKFonts}=8
\setCJKmainfont{Noto Serif CJK HK}[Ligatures=TeX]
\fi\fi\fi\fi\fi\fi\fi\fi\else
\ifnum\value{CJKLanguage}=3
\ifnum\value{CJKFonts}=1
\setCJKmainfont{Noto Sans CJK JP}[Ligatures=TeX]
\setCJKsansfont{Noto Serif CJK JP}[Ligatures=TeX]
\setCJKmonofont{Noto Sans CJK MonoCJKjp}[Ligatures=TeX]
\else\ifnum\value{CJKFonts}=2
\setCJKmainfont{Noto Serif CJK JP}[Ligatures=TeX]
\setCJKsansfont{Noto Sans CJK JP}[Ligatures=TeX]
\setCJKmonofont{Noto Sans CJK MonoCJKjp}[Ligatures=TeX]
\else\ifnum\value{CJKFonts}=3
\setCJKmainfont{Noto Sans CJK JP}[Ligatures=TeX]
\setCJKmonofont{Noto Sans CJK MonoCJKjp}[Ligatures=TeX]
\else\ifnum\value{CJKFonts}=4
\setCJKmainfont{Noto Serif CJK JP}[Ligatures=TeX]
\setCJKmonofont{Noto Sans CJK MonoCJKjp}[Ligatures=TeX]
\else\ifnum\value{CJKFonts}=5
\setCJKmainfont{Noto Sans CJK JP}[Ligatures=TeX]
\setCJKsansfont{Noto Serif CJK JP}[Ligatures=TeX]
\else\ifnum\value{CJKFonts}=6
\setCJKmainfont{Noto Serif CJK JP}[Ligatures=TeX]
\setCJKsansfont{Noto Sans CJK JP}[Ligatures=TeX]
\else\ifnum\value{CJKFonts}=7
\setCJKmainfont{Noto Sans CJK JP}[Ligatures=TeX]
\else\ifnum\value{CJKFonts}=8
\setCJKmainfont{Noto Serif CJK JP}[Ligatures=TeX]
\fi\fi\fi\fi\fi\fi\fi\fi\else
\ifnum\value{CJKLanguage}=4
\ifnum\value{CJKFonts}=1
\setCJKmainfont{Noto Sans CJK KR}[Ligatures=TeX]
\setCJKsansfont{Noto Serif CJK KR}[Ligatures=TeX]
\setCJKmonofont{Noto Sans CJK MonoCJKkr}[Ligatures=TeX]
\else\ifnum\value{CJKFonts}=2
\setCJKmainfont{Noto Serif CJK KR}[Ligatures=TeX]
\setCJKsansfont{Noto Sans CJK KR}[Ligatures=TeX]
\setCJKmonofont{Noto Sans CJK MonoCJKkr}[Ligatures=TeX]
\else\ifnum\value{CJKFonts}=3
\setCJKmainfont{Noto Sans CJK KR}[Ligatures=TeX]
\setCJKmonofont{Noto Sans CJK MonoCJKkr}[Ligatures=TeX]
\else\ifnum\value{CJKFonts}=4
\setCJKmainfont{Noto Serif CJK KR}[Ligatures=TeX]
\setCJKmonofont{Noto Sans CJK MonoCJKkr}[Ligatures=TeX]
\else\ifnum\value{CJKFonts}=5
\setCJKmainfont{Noto Sans CJK KR}[Ligatures=TeX]
\setCJKsansfont{Noto Serif CJK KR}[Ligatures=TeX]
\else\ifnum\value{CJKFonts}=6
\setCJKmainfont{Noto Serif CJK KR}[Ligatures=TeX]
\setCJKsansfont{Noto Sans CJK KR}[Ligatures=TeX]
\else\ifnum\value{CJKFonts}=7
\setCJKmainfont{Noto Sans CJK KR}[Ligatures=TeX]
\else\ifnum\value{CJKFonts}=8
\setCJKmainfont{Noto Serif CJK KR}[Ligatures=TeX]
\fi\fi\fi\fi\fi\fi\fi\fi\fi\fi\fi\fi\fi\fi
\ifnum\value{NotoCJKFamily}=0
\else\ifnum\value{NotoCJKFamily}=1
\newfontfamily\SansTC{Noto Sans CJK TC}[Ligatures=TeX]
\newfontfamily\TC{Noto Sans CJK TC}[Ligatures=TeX]
\newfontfamily\SerifTC{Noto Serif CJK TC}[Ligatures=TeX]
\newfontfamily\MonoTC{Noto Sans CJK MonoCJKtc}[Ligatures=TeX]
\newfontfamily\SansSC{Noto Sans CJK SC}[Ligatures=TeX]
\newfontfamily\SC{Noto Sans CJK SC}[Ligatures=TeX]
\newfontfamily\SerifSC{Noto Serif CJK SC}[Ligatures=TeX]
\newfontfamily\MonoSC{Noto Sans CJK MonoCJKsc}[Ligatures=TeX]
\newfontfamily\SansHK{Noto Sans CJK HK}[Ligatures=TeX]
\newfontfamily\HK{Noto Sans CJK HK}[Ligatures=TeX]
\newfontfamily\SerifHK{Noto Serif CJK HK}[Ligatures=TeX]
\newfontfamily\MonoHK{Noto Sans CJK MonoCJKhk}[Ligatures=TeX]
\newfontfamily\SansJP{Noto Sans CJK JP}[Ligatures=TeX]
\newfontfamily\JP{Noto Sans CJK JP}[Ligatures=TeX]
\newfontfamily\SerifJP{Noto Serif CJK JP}[Ligatures=TeX]
\newfontfamily\MonoJP{Noto Sans CJK MonoCJKjp}[Ligatures=TeX]
\newfontfamily\SansKR{Noto Sans CJK KR}[Ligatures=TeX]
\newfontfamily\KR{Noto Sans CJK KR}[Ligatures=TeX]
\newfontfamily\SerifKR{Noto Serif CJK KR}[Ligatures=TeX]
\newfontfamily\MonoKR{Noto Sans CJK MonoCJKkr}[Ligatures=TeX]\else\ifnum\value{NotoCJKFamily}=2
\newfontfamily\SansTC{Noto Sans CJK TC}[Ligatures=TeX]
\newfontfamily\SerifTC{Noto Serif CJK TC}[Ligatures=TeX]
\newfontfamily\TC{Noto Serif CJK TC}[Ligatures=TeX]
\newfontfamily\MonoTC{Noto Sans CJK MonoCJKtc}[Ligatures=TeX]
\newfontfamily\SansSC{Noto Sans CJK SC}[Ligatures=TeX]
\newfontfamily\SerifSC{Noto Serif CJK SC}[Ligatures=TeX]
\newfontfamily\SC{Noto Serif CJK SC}[Ligatures=TeX]
\newfontfamily\MonoSC{Noto Sans CJK MonoCJKsc}[Ligatures=TeX]
\newfontfamily\SansHK{Noto Sans CJK HK}[Ligatures=TeX]
\newfontfamily\SerifHK{Noto Serif CJK HK}[Ligatures=TeX]
\newfontfamily\HK{Noto Serif CJK HK}[Ligatures=TeX]
\newfontfamily\MonoHK{Noto Sans CJK MonoCJKhk}[Ligatures=TeX]
\newfontfamily\SansJP{Noto Sans CJK JP}[Ligatures=TeX]
\newfontfamily\SerifJP{Noto Serif CJK JP}[Ligatures=TeX]
\newfontfamily\JP{Noto Serif CJK JP}[Ligatures=TeX]
\newfontfamily\MonoJP{Noto Sans CJK MonoCJKjp}[Ligatures=TeX]
\newfontfamily\SansKR{Noto Sans CJK KR}[Ligatures=TeX]
\newfontfamily\SerifKR{Noto Serif CJK KR}[Ligatures=TeX]
\newfontfamily\KR{Noto Serif CJK KR}[Ligatures=TeX]
\newfontfamily\MonoKR{Noto Sans CJK MonoCJKkr}[Ligatures=TeX]\else\ifnum\value{NotoCJKFamily}=3
\newfontfamily\SansTC{Noto Sans CJK TC}[Ligatures=TeX]
\newfontfamily\SerifTC{Noto Serif CJK TC}[Ligatures=TeX]
\newfontfamily\MonoTC{Noto Sans CJK MonoCJKtc}[Ligatures=TeX]
\newfontfamily\SansSC{Noto Sans CJK SC}[Ligatures=TeX]
\newfontfamily\SerifSC{Noto Serif CJK SC}[Ligatures=TeX]
\newfontfamily\MonoSC{Noto Sans CJK MonoCJKsc}[Ligatures=TeX]
\newfontfamily\SansHK{Noto Sans CJK HK}[Ligatures=TeX]
\newfontfamily\SerifHK{Noto Serif CJK HK}[Ligatures=TeX]
\newfontfamily\MonoHK{Noto Sans CJK MonoCJKhk}[Ligatures=TeX]
\newfontfamily\SansJP{Noto Sans CJK JP}[Ligatures=TeX]
\newfontfamily\SerifJP{Noto Serif CJK JP}[Ligatures=TeX]
\newfontfamily\MonoJP{Noto Sans CJK MonoCJKjp}[Ligatures=TeX]
\newfontfamily\SansKR{Noto Sans CJK KR}[Ligatures=TeX]
\newfontfamily\SerifKR{Noto Serif CJK KR}[Ligatures=TeX]
\newfontfamily\MonoKR{Noto Sans CJK MonoCJKkr}[Ligatures=TeX]\else\ifnum\value{NotoCJKFamily}=4
\newfontfamily\SansTC{Noto Sans CJK TC}[Ligatures=TeX]
\newfontfamily\TC{Noto Sans CJK TC}[Ligatures=TeX]
\newfontfamily\SerifTC{Noto Serif CJK TC}[Ligatures=TeX]
\newfontfamily\SansSC{Noto Sans CJK SC}[Ligatures=TeX]
\newfontfamily\SC{Noto Sans CJK SC}[Ligatures=TeX]
\newfontfamily\SerifSC{Noto Serif CJK SC}[Ligatures=TeX]
\newfontfamily\SansHK{Noto Sans CJK HK}[Ligatures=TeX]
\newfontfamily\HK{Noto Sans CJK HK}[Ligatures=TeX]
\newfontfamily\SerifHK{Noto Serif CJK HK}[Ligatures=TeX]
\newfontfamily\SansJP{Noto Sans CJK JP}[Ligatures=TeX]
\newfontfamily\JP{Noto Sans CJK JP}[Ligatures=TeX]
\newfontfamily\SerifJP{Noto Serif CJK JP}[Ligatures=TeX]
\newfontfamily\SansKR{Noto Sans CJK KR}[Ligatures=TeX]
\newfontfamily\KR{Noto Sans CJK KR}[Ligatures=TeX]
\newfontfamily\SerifKR{Noto Serif CJK KR}[Ligatures=TeX]\else\ifnum\value{NotoCJKFamily}=5
\newfontfamily\SansTC{Noto Sans CJK TC}[Ligatures=TeX]
\newfontfamily\SerifTC{Noto Serif CJK TC}[Ligatures=TeX]
\newfontfamily\TC{Noto Serif CJK TC}[Ligatures=TeX]
\newfontfamily\SansSC{Noto Sans CJK SC}[Ligatures=TeX]
\newfontfamily\SerifSC{Noto Serif CJK SC}[Ligatures=TeX]
\newfontfamily\SC{Noto Serif CJK SC}[Ligatures=TeX]
\newfontfamily\SansHK{Noto Sans CJK HK}[Ligatures=TeX]
\newfontfamily\SerifHK{Noto Serif CJK HK}[Ligatures=TeX]
\newfontfamily\HK{Noto Serif CJK HK}[Ligatures=TeX]
\newfontfamily\SansJP{Noto Sans CJK JP}[Ligatures=TeX]
\newfontfamily\SerifJP{Noto Serif CJK JP}[Ligatures=TeX]
\newfontfamily\JP{Noto Serif CJK JP}[Ligatures=TeX]
\newfontfamily\SansKR{Noto Sans CJK KR}[Ligatures=TeX]
\newfontfamily\SerifKR{Noto Serif CJK KR}[Ligatures=TeX]
\newfontfamily\KR{Noto Serif CJK KR}[Ligatures=TeX]\else\ifnum\value{NotoCJKFamily}=6
\newfontfamily\SansTC{Noto Sans CJK TC}[Ligatures=TeX]
\newfontfamily\SerifTC{Noto Serif CJK TC}[Ligatures=TeX]
\newfontfamily\SansSC{Noto Sans CJK SC}[Ligatures=TeX]
\newfontfamily\SerifSC{Noto Serif CJK SC}[Ligatures=TeX]
\newfontfamily\SansHK{Noto Sans CJK HK}[Ligatures=TeX]
\newfontfamily\SerifHK{Noto Serif CJK HK}[Ligatures=TeX]
\newfontfamily\SansJP{Noto Sans CJK JP}[Ligatures=TeX]
\newfontfamily\SerifJP{Noto Serif CJK JP}[Ligatures=TeX]
\newfontfamily\SansKR{Noto Sans CJK KR}[Ligatures=TeX]
\newfontfamily\SerifKR{Noto Serif CJK KR}[Ligatures=TeX]
\fi\fi\fi\fi\fi\fi\fi
\else
\ifnum\value{CJKFonts}=0\else
\ifnum\value{CJKLanguage}=0
\ifnum\value{CJKFonts}=1
\setmainjfont{Noto Sans CJK TC}[Ligatures=TeX]
\setsansjfont{Noto Serif CJK TC}[Ligatures=TeX]
\setmonojfont{Noto Sans CJK MonoCJKtc}[Ligatures=TeX]
\else\ifnum\value{CJKFonts}=2
\setmainjfont{Noto Serif CJK TC}[Ligatures=TeX]
\setsansjfont{Noto Sans CJK TC}[Ligatures=TeX]
\setmonojfont{Noto Sans CJK MonoCJKtc}[Ligatures=TeX]
\else\ifnum\value{CJKFonts}=3
\setmainjfont{Noto Sans CJK TC}[Ligatures=TeX]
\setmonojfont{Noto Sans CJK MonoCJKtc}[Ligatures=TeX]
\else\ifnum\value{CJKFonts}=4
\setmainjfont{Noto Serif CJK TC}[Ligatures=TeX]
\setmonojfont{Noto Sans CJK MonoCJKtc}[Ligatures=TeX]
\else\ifnum\value{CJKFonts}=5
\setmainjfont{Noto Sans CJK TC}[Ligatures=TeX]
\setsansjfont{Noto Serif CJK TC}[Ligatures=TeX]
\else\ifnum\value{CJKFonts}=6
\setmainjfont{Noto Serif CJK TC}[Ligatures=TeX]
\setsansjfont{Noto Sans CJK TC}[Ligatures=TeX]
\else\ifnum\value{CJKFonts}=7
\setmainjfont{Noto Sans CJK TC}[Ligatures=TeX]
\else\ifnum\value{CJKFonts}=8
\setmainjfont{Noto Serif CJK TC}[Ligatures=TeX]
\fi\fi\fi\fi\fi\fi\fi\fi\else
\ifnum\value{CJKLanguage}=1
\ifnum\value{CJKFonts}=1
\setmainjfont{Noto Sans CJK SC}[Ligatures=TeX]
\setsansjfont{Noto Serif CJK SC}[Ligatures=TeX]
\setmonojfont{Noto Sans CJK MonoCJKsc}[Ligatures=TeX]
\else\ifnum\value{CJKFonts}=2
\setmainjfont{Noto Serif CJK SC}[Ligatures=TeX]
\setsansjfont{Noto Sans CJK SC}[Ligatures=TeX]
\setmonojfont{Noto Sans CJK MonoCJKsc}[Ligatures=TeX]
\else\ifnum\value{CJKFonts}=3
\setmainjfont{Noto Sans CJK SC}[Ligatures=TeX]
\setmonojfont{Noto Sans CJK MonoCJKsc}[Ligatures=TeX]
\else\ifnum\value{CJKFonts}=4
\setmainjfont{Noto Serif CJK SC}[Ligatures=TeX]
\setmonojfont{Noto Sans CJK MonoCJKsc}[Ligatures=TeX]
\else\ifnum\value{CJKFonts}=5
\setmainjfont{Noto Sans CJK SC}[Ligatures=TeX]
\setsansjfont{Noto Serif CJK SC}[Ligatures=TeX]
\else\ifnum\value{CJKFonts}=6
\setmainjfont{Noto Serif CJK SC}[Ligatures=TeX]
\setsansjfont{Noto Sans CJK SC}[Ligatures=TeX]
\else\ifnum\value{CJKFonts}=7
\setmainjfont{Noto Sans CJK SC}[Ligatures=TeX]
\else\ifnum\value{CJKFonts}=8
\setmainjfont{Noto Serif CJK SC}[Ligatures=TeX]
\fi\fi\fi\fi\fi\fi\fi\fi\else
\ifnum\value{CJKLanguage}=2
\ifnum\value{CJKFonts}=1
\setmainjfont{Noto Sans CJK HK}[Ligatures=TeX]
\setsansjfont{Noto Serif CJK HK}[Ligatures=TeX]
\setmonojfont{Noto Sans CJK MonoCJKhk}[Ligatures=TeX]
\else\ifnum\value{CJKFonts}=2
\setmainjfont{Noto Serif CJK HK}[Ligatures=TeX]
\setsansjfont{Noto Sans CJK HK}[Ligatures=TeX]
\setmonojfont{Noto Sans CJK MonoCJKhk}[Ligatures=TeX]
\else\ifnum\value{CJKFonts}=3
\setmainjfont{Noto Sans CJK HK}[Ligatures=TeX]
\setmonojfont{Noto Sans CJK MonoCJKhk}[Ligatures=TeX]
\else\ifnum\value{CJKFonts}=4
\setmainjfont{Noto Serif CJK HK}[Ligatures=TeX]
\setmonojfont{Noto Sans CJK MonoCJKhk}[Ligatures=TeX]
\else\ifnum\value{CJKFonts}=5
\setmainjfont{Noto Sans CJK HK}[Ligatures=TeX]
\setsansjfont{Noto Serif CJK HK}[Ligatures=TeX]
\else\ifnum\value{CJKFonts}=6
\setmainjfont{Noto Serif CJK HK}[Ligatures=TeX]
\setsansjfont{Noto Sans CJK HK}[Ligatures=TeX]
\else\ifnum\value{CJKFonts}=7
\setmainjfont{Noto Sans CJK HK}[Ligatures=TeX]
\else\ifnum\value{CJKFonts}=8
\setmainjfont{Noto Serif CJK HK}[Ligatures=TeX]
\fi\fi\fi\fi\fi\fi\fi\fi\else
\ifnum\value{CJKLanguage}=3
\ifnum\value{CJKFonts}=1
\setmainjfont{Noto Sans CJK JP}[Ligatures=TeX]
\setsansjfont{Noto Serif CJK JP}[Ligatures=TeX]
\setmonojfont{Noto Sans CJK MonoCJKjp}[Ligatures=TeX]
\else\ifnum\value{CJKFonts}=2
\setmainjfont{Noto Serif CJK JP}[Ligatures=TeX]
\setsansjfont{Noto Sans CJK JP}[Ligatures=TeX]
\setmonojfont{Noto Sans CJK MonoCJKjp}[Ligatures=TeX]
\else\ifnum\value{CJKFonts}=3
\setmainjfont{Noto Sans CJK JP}[Ligatures=TeX]
\setmonojfont{Noto Sans CJK MonoCJKjp}[Ligatures=TeX]
\else\ifnum\value{CJKFonts}=4
\setmainjfont{Noto Serif CJK JP}[Ligatures=TeX]
\setmonojfont{Noto Sans CJK MonoCJKjp}[Ligatures=TeX]
\else\ifnum\value{CJKFonts}=5
\setmainjfont{Noto Sans CJK JP}[Ligatures=TeX]
\setsansjfont{Noto Serif CJK JP}[Ligatures=TeX]
\else\ifnum\value{CJKFonts}=6
\setmainjfont{Noto Serif CJK JP}[Ligatures=TeX]
\setsansjfont{Noto Sans CJK JP}[Ligatures=TeX]
\else\ifnum\value{CJKFonts}=7
\setmainjfont{Noto Sans CJK JP}[Ligatures=TeX]
\else\ifnum\value{CJKFonts}=8
\setmainjfont{Noto Serif CJK JP}[Ligatures=TeX]
\fi\fi\fi\fi\fi\fi\fi\fi\else
\ifnum\value{CJKLanguage}=4
\ifnum\value{CJKFonts}=1
\setmainjfont{Noto Sans CJK KR}[Ligatures=TeX]
\setsansjfont{Noto Serif CJK KR}[Ligatures=TeX]
\setmonojfont{Noto Sans CJK MonoCJKkr}[Ligatures=TeX]
\else\ifnum\value{CJKFonts}=2
\setmainjfont{Noto Serif CJK KR}[Ligatures=TeX]
\setsansjfont{Noto Sans CJK KR}[Ligatures=TeX]
\setmonojfont{Noto Sans CJK MonoCJKkr}[Ligatures=TeX]
\else\ifnum\value{CJKFonts}=3
\setmainjfont{Noto Sans CJK KR}[Ligatures=TeX]
\setmonojfont{Noto Sans CJK MonoCJKkr}[Ligatures=TeX]
\else\ifnum\value{CJKFonts}=4
\setmainjfont{Noto Serif CJK KR}[Ligatures=TeX]
\setmonojfont{Noto Sans CJK MonoCJKkr}[Ligatures=TeX]
\else\ifnum\value{CJKFonts}=5
\setmainjfont{Noto Sans CJK KR}[Ligatures=TeX]
\setsansjfont{Noto Serif CJK KR}[Ligatures=TeX]
\else\ifnum\value{CJKFonts}=6
\setmainjfont{Noto Serif CJK KR}[Ligatures=TeX]
\setsansjfont{Noto Sans CJK KR}[Ligatures=TeX]
\else\ifnum\value{CJKFonts}=7
\setmainjfont{Noto Sans CJK KR}[Ligatures=TeX]
\else\ifnum\value{CJKFonts}=8
\setmainjfont{Noto Serif CJK KR}[Ligatures=TeX]
\fi\fi\fi\fi\fi\fi\fi\fi\fi\fi\fi\fi\fi\fi
\ifnum\value{NotoCJKFamily}=0
\else\ifnum\value{NotoCJKFamily}=1
\newfontfamily\SansTC{Noto Sans CJK TC}[Ligatures=TeX]
\newfontfamily\TC{Noto Sans CJK TC}[Ligatures=TeX]
\newfontfamily\SerifTC{Noto Serif CJK TC}[Ligatures=TeX]
\newfontfamily\MonoTC{Noto Sans CJK MonoCJKtc}[Ligatures=TeX]
\newfontfamily\SansSC{Noto Sans CJK SC}[Ligatures=TeX]
\newfontfamily\SC{Noto Sans CJK SC}[Ligatures=TeX]
\newfontfamily\SerifSC{Noto Serif CJK SC}[Ligatures=TeX]
\newfontfamily\MonoSC{Noto Sans CJK MonoCJKsc}[Ligatures=TeX]
\newfontfamily\SansHK{Noto Sans CJK HK}[Ligatures=TeX]
\newfontfamily\HK{Noto Sans CJK HK}[Ligatures=TeX]
\newfontfamily\SerifHK{Noto Serif CJK HK}[Ligatures=TeX]
\newfontfamily\MonoHK{Noto Sans CJK MonoCJKhk}[Ligatures=TeX]
\newfontfamily\SansJP{Noto Sans CJK JP}[Ligatures=TeX]
\newfontfamily\JP{Noto Sans CJK JP}[Ligatures=TeX]
\newfontfamily\SerifJP{Noto Serif CJK JP}[Ligatures=TeX]
\newfontfamily\MonoJP{Noto Sans CJK MonoCJKjp}[Ligatures=TeX]
\newfontfamily\SansKR{Noto Sans CJK KR}[Ligatures=TeX]
\newfontfamily\KR{Noto Sans CJK KR}[Ligatures=TeX]
\newfontfamily\SerifKR{Noto Serif CJK KR}[Ligatures=TeX]
\newfontfamily\MonoKR{Noto Sans CJK MonoCJKkr}[Ligatures=TeX]\else\ifnum\value{NotoCJKFamily}=2
\newfontfamily\SansTC{Noto Sans CJK TC}[Ligatures=TeX]
\newfontfamily\SerifTC{Noto Serif CJK TC}[Ligatures=TeX]
\newfontfamily\TC{Noto Serif CJK TC}[Ligatures=TeX]
\newfontfamily\MonoTC{Noto Sans CJK MonoCJKtc}[Ligatures=TeX]
\newfontfamily\SansSC{Noto Sans CJK SC}[Ligatures=TeX]
\newfontfamily\SerifSC{Noto Serif CJK SC}[Ligatures=TeX]
\newfontfamily\SC{Noto Serif CJK SC}[Ligatures=TeX]
\newfontfamily\MonoSC{Noto Sans CJK MonoCJKsc}[Ligatures=TeX]
\newfontfamily\SansHK{Noto Sans CJK HK}[Ligatures=TeX]
\newfontfamily\SerifHK{Noto Serif CJK HK}[Ligatures=TeX]
\newfontfamily\HK{Noto Serif CJK HK}[Ligatures=TeX]
\newfontfamily\MonoHK{Noto Sans CJK MonoCJKhk}[Ligatures=TeX]
\newfontfamily\SansJP{Noto Sans CJK JP}[Ligatures=TeX]
\newfontfamily\SerifJP{Noto Serif CJK JP}[Ligatures=TeX]
\newfontfamily\JP{Noto Serif CJK JP}[Ligatures=TeX]
\newfontfamily\MonoJP{Noto Sans CJK MonoCJKjp}[Ligatures=TeX]
\newfontfamily\SansKR{Noto Sans CJK KR}[Ligatures=TeX]
\newfontfamily\SerifKR{Noto Serif CJK KR}[Ligatures=TeX]
\newfontfamily\KR{Noto Serif CJK KR}[Ligatures=TeX]
\newfontfamily\MonoKR{Noto Sans CJK MonoCJKkr}[Ligatures=TeX]\else\ifnum\value{NotoCJKFamily}=3
\newfontfamily\SansTC{Noto Sans CJK TC}[Ligatures=TeX]
\newfontfamily\SerifTC{Noto Serif CJK TC}[Ligatures=TeX]
\newfontfamily\MonoTC{Noto Sans CJK MonoCJKtc}[Ligatures=TeX]
\newfontfamily\SansSC{Noto Sans CJK SC}[Ligatures=TeX]
\newfontfamily\SerifSC{Noto Serif CJK SC}[Ligatures=TeX]
\newfontfamily\MonoSC{Noto Sans CJK MonoCJKsc}[Ligatures=TeX]
\newfontfamily\SansHK{Noto Sans CJK HK}[Ligatures=TeX]
\newfontfamily\SerifHK{Noto Serif CJK HK}[Ligatures=TeX]
\newfontfamily\MonoHK{Noto Sans CJK MonoCJKhk}[Ligatures=TeX]
\newfontfamily\SansJP{Noto Sans CJK JP}[Ligatures=TeX]
\newfontfamily\SerifJP{Noto Serif CJK JP}[Ligatures=TeX]
\newfontfamily\MonoJP{Noto Sans CJK MonoCJKjp}[Ligatures=TeX]
\newfontfamily\SansKR{Noto Sans CJK KR}[Ligatures=TeX]
\newfontfamily\SerifKR{Noto Serif CJK KR}[Ligatures=TeX]
\newfontfamily\MonoKR{Noto Sans CJK MonoCJKkr}[Ligatures=TeX]\else\ifnum\value{NotoCJKFamily}=4
\newfontfamily\SansTC{Noto Sans CJK TC}[Ligatures=TeX]
\newfontfamily\TC{Noto Sans CJK TC}[Ligatures=TeX]
\newfontfamily\SerifTC{Noto Serif CJK TC}[Ligatures=TeX]
\newfontfamily\SansSC{Noto Sans CJK SC}[Ligatures=TeX]
\newfontfamily\SC{Noto Sans CJK SC}[Ligatures=TeX]
\newfontfamily\SerifSC{Noto Serif CJK SC}[Ligatures=TeX]
\newfontfamily\SansHK{Noto Sans CJK HK}[Ligatures=TeX]
\newfontfamily\HK{Noto Sans CJK HK}[Ligatures=TeX]
\newfontfamily\SerifHK{Noto Serif CJK HK}[Ligatures=TeX]
\newfontfamily\SansJP{Noto Sans CJK JP}[Ligatures=TeX]
\newfontfamily\JP{Noto Sans CJK JP}[Ligatures=TeX]
\newfontfamily\SerifJP{Noto Serif CJK JP}[Ligatures=TeX]
\newfontfamily\SansKR{Noto Sans CJK KR}[Ligatures=TeX]
\newfontfamily\KR{Noto Sans CJK KR}[Ligatures=TeX]
\newfontfamily\SerifKR{Noto Serif CJK KR}[Ligatures=TeX]\else\ifnum\value{NotoCJKFamily}=5
\newfontfamily\SansTC{Noto Sans CJK TC}[Ligatures=TeX]
\newfontfamily\SerifTC{Noto Serif CJK TC}[Ligatures=TeX]
\newfontfamily\TC{Noto Serif CJK TC}[Ligatures=TeX]
\newfontfamily\SansSC{Noto Sans CJK SC}[Ligatures=TeX]
\newfontfamily\SerifSC{Noto Serif CJK SC}[Ligatures=TeX]
\newfontfamily\SC{Noto Serif CJK SC}[Ligatures=TeX]
\newfontfamily\SansHK{Noto Sans CJK HK}[Ligatures=TeX]
\newfontfamily\SerifHK{Noto Serif CJK HK}[Ligatures=TeX]
\newfontfamily\HK{Noto Serif CJK HK}[Ligatures=TeX]
\newfontfamily\SansJP{Noto Sans CJK JP}[Ligatures=TeX]
\newfontfamily\SerifJP{Noto Serif CJK JP}[Ligatures=TeX]
\newfontfamily\JP{Noto Serif CJK JP}[Ligatures=TeX]
\newfontfamily\SansKR{Noto Sans CJK KR}[Ligatures=TeX]
\newfontfamily\SerifKR{Noto Serif CJK KR}[Ligatures=TeX]
\newfontfamily\KR{Noto Serif CJK KR}[Ligatures=TeX]\else\ifnum\value{NotoCJKFamily}=6
\newfontfamily\SansTC{Noto Sans CJK TC}[Ligatures=TeX]
\newfontfamily\SerifTC{Noto Serif CJK TC}[Ligatures=TeX]
\newfontfamily\SansSC{Noto Sans CJK SC}[Ligatures=TeX]
\newfontfamily\SerifSC{Noto Serif CJK SC}[Ligatures=TeX]
\newfontfamily\SansHK{Noto Sans CJK HK}[Ligatures=TeX]
\newfontfamily\SerifHK{Noto Serif CJK HK}[Ligatures=TeX]
\newfontfamily\SansJP{Noto Sans CJK JP}[Ligatures=TeX]
\newfontfamily\SerifJP{Noto Serif CJK JP}[Ligatures=TeX]
\newfontfamily\SansKR{Noto Sans CJK KR}[Ligatures=TeX]
\newfontfamily\SerifKR{Noto Serif CJK KR}[Ligatures=TeX]
\fi\fi\fi\fi\fi\fi\fi
\fi\fi
\righthyphenmin=4
\lefthyphenmin=4
\pretolerance=5000
\tolerance=9000
\emergencystretch=1em
\setlength{\parindent}{0pt}
\everymath{\displaystyle}
\everydisplay{\displaystyle}
\captionsetup{labelformat=empty}
\renewcommand{\maketitle}{
\begin{titlepage}
\begin{center}
\vspace*{\fill}
{\huge \bfseries \thetitle\par}
\vskip 1.5em
{\Large \theauthor\par}
\vskip 1em
{\large \thedate\par}
\vspace*{\fill}
\end{center}
\end{titlepage}
}
\lstset{
basicstyle=\ttfamily\scriptsize,
keywordstyle=\color{blue},
stringstyle=\color{red},
commentstyle=\color{green!50!black},
numberstyle=\small\color{gray},
numbers=left,
stepnumber=1,
numbersep=5pt,
backgroundcolor=\color{white},
showspaces=false,
showstringspaces=false,
showtabs=false,
frame=single,
rulecolor=\color{black},
tabsize=4,
captionpos=b,
breaklines=true,
breakatwhitespace=false
}
\makeatletter\csvset{autotabularcenter/.style={file=#1,after head=\csv@pretable\begin{tabular}{|*{\csv@columncount}{c|}}\csv@tablehead,table head=\hline\csvlinetotablerow\\\hline,late after line=\\,lable foot=\\\hline,late after last line=\csv@tablefoot\end{tabular}\csv@posttable,command=\csvlinetotablerow},autobooktabularcenter/.style={file=#1,after head=\csv@pretable\begin{tabular}{*{\csv@columncount}{c}}\csv@tablehead,table head=\toprule\csvlinetotablerow\\\midrule,late after line=\\,table foot=\\\bottomrule,late after last line=\csv@tablefoot\end{tabular}\csv@posttable,command=\csvlinetotablerow}}\makeatother
\newcommand{\csvautotabularcenter}[2][]{\csvloop{autotabularcenter={#2},#1}}
\newcommand{\csvautobooktabularcenter}[2][]{\csvloop{autobooktabularcenter={#2},#1}}
\newtheorem{theorem}{Theorem}
\newtheorem{lemma}{Lemma}
\newcolumntype{C}[1]{>{\centering\arraybackslash}p{#1}}
\newcolumntype{L}[1]{>{\raggedright\arraybackslash}p{#1}}
\newcolumntype{R}[1]{>{\raggedleft\arraybackslash}p{#1}}
\newcolumntype{m}{>{$\displaystyle}c<{$}}
\newcolumntype{M}[1]{>{$\displaystyle}p{#1}<{$}}
\setlength{\LTleft}{\fill}
\setlength{\LTright}{\fill}
\makeatletter
\tikzset{dot diameter/.store in=\dot@diameter,dot diameter=1pt,dot spacing/.store in=\dot@spacing,dot spacing=5.5pt,dots/.style={line width=\dot@diameter,line cap=round,dash pattern=on 0pt off \dot@spacing}}
\newcommand{\overleftrightarrowsmall}{\mathpalette{\overarrowsmall@\leftrightarrowfill@}}
\newcommand{\overrightarrowsmall}{\mathpalette{\overarrowsmall@\rightarrowfill@}}
\newcommand{\overleftarrowsmall}{\mathpalette{\overarrowsmall@\leftarrowfill@}}
\newcommand{\overarrowsmall@}[3]{\vbox{\ialign{##\crcr#1{\smaller@style{#2}}\crcr\noalign{\nointerlineskip}$\m@th\hfil#2#3\hfil$\crcr}}}\def\smaller@style#1{\ifx#1\displaystyle\scriptstyle\else\ifx#1\textstyle\scriptstyle\else\scriptscriptstyle\fi\fi}
\DeclareFontFamily{U}{tipa}{}
\DeclareFontShape{U}{tipa}{m}{n}{<->tipa10}{}
\DeclareDocumentCommand\arc@char{}{{\usefont{U}{tipa}{m}{n}\symbol{62}}}
\DeclareDocumentCommand\arc{m}{\mathpalette\arc@arc{#1}}
\DeclareDocumentCommand\arc@arc{m m}{\sbox0{$\m@th#1#2$}\vbox{\hbox{\resizebox{\wd0}{\height}{\arc@char}}\nointerlineskip\box0}}
\makeatother
\let\oldunderline\underline
\let\oldoverline\overline
\newcommand{\overlinesmall}[1]{\mskip.5\thinmuskip\oldoverline{\mskip-.5\thinmuskip{#1}\mskip-.5\thinmuskip}\mskip.5\thinmuskip}
\newcommand{\underlinesmall}[1]{\mskip.5\thinmuskip\oldunderline{\mskip-.5\thinmuskip{#1}\mskip-.5\thinmuskip}\mskip.5\thinmuskip}
\DeclareDocumentCommand\olra{}{\overleftrightarrowsmall}
\DeclareDocumentCommand\oldoverleftrightarrow{}{\overleftrightarrow}
\DeclareDocumentCommand\overleftrightarrow{}{\overleftrightarrowsmall}
\DeclareDocumentCommand\ora{}{\overrightarrowsmall}
\DeclareDocumentCommand\oldoverrightarrow{}{\overrightarrow}
\DeclareDocumentCommand\overrightarrow{}{\overrightarrowsmall}
\DeclareDocumentCommand\ola{}{\overleftarrowsmall}
\DeclareDocumentCommand\oldoverleftarrow{}{\overleftarrow}
\DeclareDocumentCommand\overleftarrow{}{\overleftarrowsmall}
\DeclareDocumentCommand\ol{}{\overlinesmall}
\DeclareDocumentCommand\overline{}{\overlinesmall}
\DeclareDocumentCommand\ul{}{\underlinesmall}
\DeclareDocumentCommand\underline{}{\underlinesmall}
\let\amp\&
\let\CE\ce
\let\cd\cdot
\let\cds\cdots
\let\chm\chemfig
\let\clp\clearpage
\let\cpt\caption
\let\ctr\centering
\let\dds\ddots
\let\dgb\diagbox
\let\FB\FloatBarrier
\let\fr\frac
\let\hf\hfill
\let\hl\hline
\let\hs\hspace
\let\icg\includegraphics
\let\icp\includepdf
\let\ift\infty
\let\ip\input
\let\lds\ldots
\let\lh\lineheight
\let\lw\linewidth
\let\mb\mathbf
\let\mbb\mathbb
\let\mbbi\mathbbit
\let\mbbit\mathbbit
\let\mbc\mathbfcal
\let\mbf\mathbf
\let\mbfcal\mathbfcal
\let\mbffrak\mathbffrak
\let\mbfit\mathbfit
\let\mbfk\mathbffrak
\let\mbfscr\mathbfscr
\let\mbi\mathbfit
\let\mbs\mathbfscr
\let\mc\mathcal
\let\mcal\mathcal
\let\mucol\multicolumn
\let\murow\multirow
\let\mds\mathds
\let\mfk\mathfrak
\let\mfrak\mathfrak
\let\mi\mathit
\let\mit\mathit
\let\mkc\makecell
\let\mnm\mathnormal
\let\mr\mathrm
\let\ms\mathscr
\let\mscr\mathscr
\let\msf\mathsf
\let\mtt\mathtt
\let\mup\mathup
\let\Ns\Needspace
\let\npgb\nopagebreak
\let\opn\operatorname
\let\pc\%
\let\pgb\pagebreak
\let\pk\textperthousand
\let\prp\propto
\let\rl\relax
\let\rsb\raisebox
\let\rtb\rotatebox
\let\sbs\subset
\let\sbse\subseteq
\let\smn\setminus
\let\sps\supset
\let\spse\supseteq
\let\sq\sqrt
\let\Tb\textbf
\let\Textbf\textbf
\let\Textheight\textheight
\let\Textit\textit
\let\Textnormal\textnormal
\let\Textrm\textrm
\let\Textsc\textsc
\let\Textsf\textsf
\let\Texttt\texttt
\let\Textup\textup
\let\Textwidth\textwidth
\let\Text\text
\let\Th\textheight
\let\Ti\textit
\let\Tnm\textnormal
\let\Tr\textrm
\let\Tsc\textsc
\let\Tsf\textsf
\let\Ttt\texttt
\let\Tup\textup
\let\Tw\textwidth
\let\Tx\text
\let\tb\textbf
\let\tc\textcircled
\let\tcbox\tcolorbox
\let\texpdf\texorpdfstring
\let\tfr\tfrac
\let\th\textheight
\let\ti\textit
\let\tnm\textnormal
\let\tr\textrm
\let\tsc\textsc
\let\tsf\textsf
\let\ttt\texttt
\let\tup\textup
\let\tw\textwidth
\let\tx\text
\let\vds\vdots
\let\emset\varnothing
\let\vf\vfill
\let\vpht\vphantom
\let\vs\vspace
\DeclarePairedDelimiter{\ceil}{\lceil}{\rceil}
\DeclarePairedDelimiter{\floor}{\lfloor}{\rfloor}
\DeclarePairedDelimiter{\brangle}{\langle}{\rangle}
\DeclarePairedDelimiter{\abs}{\lvert}{\rvert}
\DeclarePairedDelimiter{\norm}{\lVert}{\rVert}
\DeclarePairedDelimiter{\brcurl}{\{}{\}}
\DeclareDocumentCommand\GS{}{(g)}
\DeclareDocumentCommand\LQ{}{(l)}
\DeclareDocumentCommand\SL{}{(s)}
\DeclareDocumentCommand\AQ{}{(aq)}
\DeclareDocumentCommand{\X}{}{\ensuremath{\times}}
\DeclareDocumentCommand{\bBNM}{}{\begin{BNiceMatrix}}
\DeclareDocumentCommand{\bBmt}{}{\begin{Bmatrix}}
\DeclareDocumentCommand{\bNM}{}{\begin{NiceMatrix}}
\DeclareDocumentCommand{\bVNM}{}{\begin{VNiceMatrix}}
\DeclareDocumentCommand{\bVmt}{}{\begin{Vmatrix}}
\DeclareDocumentCommand{\ba}{}{\begin{aligned}}
\DeclareDocumentCommand{\bbNM}{}{\begin{bNiceMatrix}}
\DeclareDocumentCommand{\bbmt}{}{\begin{bmatrix}}
\DeclareDocumentCommand{\bcs}{}{\begin{cases}}
\DeclareDocumentCommand{\bcttbrn}{}{\begin{center}\begin{figure}[h]\centering\begin{tabular}}
\DeclareDocumentCommand{\bcttbr}{}{\begin{center}\begin{figure}[H]\centering\begin{tabular}}
\DeclareDocumentCommand{\bctfn}{}{\begin{center}\begin{figure}[h]\centering}
\DeclareDocumentCommand{\bctf}{}{\begin{center}\begin{figure}[H]\centering}
\DeclareDocumentCommand{\bcttn}{}{\begin{center}\begin{table}[h]\centering}
\DeclareDocumentCommand{\bctt}{}{\begin{center}\begin{table}[H]\centering}
\DeclareDocumentCommand{\bct}{}{\begin{center}}
\DeclareDocumentCommand{\ben}{}{\begin{enumerate}}
\DeclareDocumentCommand{\beq}{}{\begin{equation}}
\DeclareDocumentCommand{\bfH}{}{\begin{figure}[H]}
\DeclareDocumentCommand{\bf}{}{\begin{figure}}
\DeclareDocumentCommand{\bg}{m}{\begin{#1}}
\DeclareDocumentCommand{\bit}{}{\begin{itemize}}
\DeclareDocumentCommand{\blt}{}{\begin{longtable}}
\DeclareDocumentCommand{\bmac}{}{\[\begin{aligned}\begin{cases}}
\DeclareDocumentCommand{\bma}{}{\[\begin{aligned}}
\DeclareDocumentCommand{\bmc}{}{\[\begin{cases}}
\DeclareDocumentCommand{\bmt}{}{\begin{matrix}}
\DeclareDocumentCommand{\bm}{}{\[}
\DeclareDocumentCommand{\bpNM}{}{\begin{pNiceMatrix}}
\DeclareDocumentCommand{\bpmt}{}{\begin{pmatrix}}
\DeclareDocumentCommand{\bpr}{}{\begin{proof}}
\DeclareDocumentCommand{\bs}{m}{\boldsymbol{#1}}
\DeclareDocumentCommand{\bsh}{}{\relax\ifmmode\backslash\else\textbackslash\fi}
\DeclareDocumentCommand{\btH}{}{\begin{table}[H]}
\DeclareDocumentCommand{\btbr}{}{\begin{tabular}}
\DeclareDocumentCommand{\btk}{}{\begin{tikzpicture}}
\DeclareDocumentCommand{\bt}{}{\begin{table}}
\DeclareDocumentCommand{\bvNM}{}{\begin{vNiceMatrix}}
\DeclareDocumentCommand{\bvmt}{}{\begin{vmatrix}}
\DeclareDocumentCommand{\bqn}{}{\begin{quotation}}
\DeclareDocumentCommand{\bqt}{}{\begin{quote}}
\DeclareDocumentCommand{\cktee}{m o}{\hspace{1em}\IfNoValueTF{#2}{\tikz[circuit ee IEC] \draw node [#1,anchor=west] {};}{\tikz[circuit ee IEC] \draw node [#1,label=#2,anchor=west] {};}\hspace{1em}}
\DeclareDocumentCommand{\cktiec}{m m o}{\hspace{1em}\IfNoValueTF{#3}{\tikz[circuit logic IEC] \draw node [#1,inputs=#2,anchor=west,name=s] {};}{\tikz[circuit logic IEC] \draw node [#1,inputs=#2,label=#3,anchor=west,name=s] {};}\hspace{1em}}
\DeclareDocumentCommand{\cktus}{m m o}{\hspace{1em}\IfNoValueTF{#3}{\tikz[circuit logic US] \draw node [#1,inputs=#2,anchor=west,name=s] {};}{\tikz[circuit logic US] \draw node [#1,inputs=#2,label=#3,anchor=west,name=s] {};}\hspace{1em}}
\DeclareDocumentCommand{\ctgn}{o m o}{\begin{center}\begin{figure}[h]\centering\IfNoValueTF{#3}{\IfNoValueTF{#2}{\includegraphics{#1}}{\includegraphics[#1]{#2}}}{\includegraphics[#1]{#2}\caption{#3}}\end{figure}\end{center}}
\DeclareDocumentCommand{\ctg}{o m o}{\begin{center}\begin{figure}[H]\centering\IfNoValueTF{#3}{\IfNoValueTF{#2}{\includegraphics{#1}}{\includegraphics[#1]{#2}}}{\includegraphics[#1]{#2}\caption{#3}}\end{figure}\FloatBarrier\end{center}}
\DeclareDocumentCommand{\df}{m m o}{\IfNoValueTF{#3}{\frac{\mathrm{d}#1}{\mathrm{d}#2}}{\frac{\mathrm{d}^{#1}#2}{\mathrm{d}#3^{#1}}}}
\DeclareDocumentCommand{\dpre}{m m o}{\IfNoValueTF{#3}{\frac{\mathrm{d}}{\mathrm{d}#2}\left(#1\right)}{\frac{\mathrm{d}^{#1}}{\mathrm{d}#3^{#1}}\left(#2\right)}}
\DeclareDocumentCommand{\eBNM}{}{\end{BNiceMatrix}}
\DeclareDocumentCommand{\eBmt}{}{\end{Bmatrix}}
\DeclareDocumentCommand{\eNM}{}{\end{NiceMatrix}}
\DeclareDocumentCommand{\eVNM}{}{\end{VNiceMatrix}}
\DeclareDocumentCommand{\eVmt}{}{\end{Vmatrix}}
\DeclareDocumentCommand{\eam}{}{\end{aligned}\]}
\DeclareDocumentCommand{\ea}{}{\end{aligned}}
\DeclareDocumentCommand{\ebNM}{}{\end{bNiceMatrix}}
\DeclareDocumentCommand{\ebmt}{}{\end{bmatrix}}
\DeclareDocumentCommand{\ecam}{}{\end{cases}\end{aligned}\]}
\DeclareDocumentCommand{\ecm}{}{\end{cases}\]}
\DeclareDocumentCommand{\ecs}{}{\end{cases}}
\DeclareDocumentCommand{\ect}{}{\end{center}}
\DeclareDocumentCommand{\ed}{m}{\end{#1}}
\DeclareDocumentCommand{\eit}{}{\end{itemize}}
\DeclareDocumentCommand{\een}{}{\end{enumerate}}
\DeclareDocumentCommand{\eeq}{}{\end{equation}}
\DeclareDocumentCommand{\efB}{}{\end{figure}\FloatBarrier}
\DeclareDocumentCommand{\efctn}{}{\end{figure}\end{center}}
\DeclareDocumentCommand{\efct}{}{\end{figure}\FloatBarrier\end{center}}
\DeclareDocumentCommand{\ef}{}{\end{figure}}
\DeclareDocumentCommand{\elt}{}{\end{longtable}}
\DeclareDocumentCommand{\emt}{}{\end{matirx}}
\DeclareDocumentCommand{\em}{}{\]}
\DeclareDocumentCommand{\epNM}{}{\end{pNiceMatrix}}
\DeclareDocumentCommand{\epmt}{}{\end{pmatrix}}
\DeclareDocumentCommand{\epr}{}{\end{proof}}
\DeclareDocumentCommand{\eq}{}{=}
\DeclareDocumentCommand{\eqn}{}{\end{quotation}}
\DeclareDocumentCommand{\eqt}{}{\end{quote}}
\DeclareDocumentCommand{\etB}{}{\end{table}\FloatBarrier}
\DeclareDocumentCommand{\etbr}{}{\end{tabular}}
\DeclareDocumentCommand{\etbrctn}{}{\end{tabular}\end{figure}\end{center}}
\DeclareDocumentCommand{\etbrct}{}{\end{tabular}\end{figure}\FloatBarrier\end{center}}
\DeclareDocumentCommand{\etctn}{}{\end{table}\end{center}}
\DeclareDocumentCommand{\etct}{}{\end{table}\FloatBarrier\end{center}}
\DeclareDocumentCommand{\etk}{}{\end{tikzpicture}}
\DeclareDocumentCommand{\et}{}{\end{table}}
\DeclareDocumentCommand{\evNM}{}{\end{vNiceMatrix}}
\DeclareDocumentCommand{\evmt}{}{\end{vmatrix}}
\DeclareDocumentCommand{\gt}{}{>}
\DeclareDocumentCommand{\INF}{}{\infty}
\DeclareDocumentCommand{\itg}{m m o o}{\IfNoValueTF{#4}{\IfNoValueTF{#3}{\int #1\,\mathrm{d}{#2}}{\int_0^{#1}#2\,\mathrm{d}{#3}}}{\int_{#1}^{#2} #3\,\mathrm{d}{#4}}}
\DeclareDocumentCommand{\lt}{}{<}
\DeclareDocumentCommand{\mdp}{O{}}{\[#1\]}
\DeclareDocumentCommand{\mil}{O{}}{\(#1\)}
\DeclareDocumentCommand{\mn}{}{-}
\DeclareDocumentCommand{\NINF}{}{-\infty}
\DeclareDocumentCommand{\nthm}{}{\Needspace{1\textheight-1em}}
\DeclareDocumentCommand{\nth}{o}{\IfNoValueTF{#1}{\Needspace{1\textheight}}{\Needspace{#1\textheight}}}
\DeclareDocumentCommand{\n}{}{\\}
\DeclareDocumentCommand{\pf}{m m o}{\IfNoValueTF{#3}{\frac{\partial #1}{\partial #2}}{\frac{\partial^{#1} #2}{\partial #3^{#1}}}}
\DeclareDocumentCommand{\pht}{m}{\phantom{#1}}
\DeclareDocumentCommand{\pl}{}{+}
\DeclareDocumentCommand{\ppre}{m m o}{\IfNoValueTF{#3}{\partial #2 \left( #1 \right)}{\frac{\partial^{#1}}{\partial #3^{#1}} \left( #2 \right)}}
\DeclareDocumentCommand{\sh}{}{/}
\DeclareDocumentCommand{\ts}{o}{\IfNoValueTF{#1}{\text{\protect\ }}{\text{\protect\foreach\i in{1,2,...,#1}{\protect\ }}}}
\DeclareDocumentCommand{\ttimes}{}{\(\times\)}
\DeclareDocumentCommand{\vsf}{}{\vspace*{\fill}}
\titleclass{\subsubparagraph}{straight}[\subparagraph]
\newcounter{subsubparagraph}[subparagraph]
\titlespacing*{\subsubparagraph}{\parindent}{3.25ex plus 1ex minus .2ex}{1em}
\titleformat{\section}[block]{\Large\bfseries}{\thesection}{1em}{}
\titleformat{\subsection}{\large\bfseries}{\thesubsection}{1em}{}
\titleformat{\subsubsection}{\normalsize\bfseries}{\thesubsubsection}{1em}{}
\titleformat{\paragraph}{\normalsize\bfseries}{\theparagraph}{1em}{}
\titleformat{\subparagraph}{\normalsize\bfseries}{\thesubparagraph}{1em}{}
\titleformat{\subsubparagraph}{\normalsize\bfseries}{\thesubsubparagraph}{1em}{}
\DeclareDocumentCommand{\sct}{s o}{\IfBooleanTF{#1}{\IfNoValueTF{#2}{\section*}{\section*{#2}}}{\IfNoValueTF{#2}{\section}{\section{#2}}}}
\DeclareDocumentCommand{\ssc}{s o}{\IfBooleanTF{#1}{\IfNoValueTF{#2}{\subsection*}{\subsection*{#2}}}{\IfNoValueTF{#2}{\subsection}{\subsection{#2}}}}
\DeclareDocumentCommand{\sssc}{s o}{\IfBooleanTF{#1}{\IfNoValueTF{#2}{\subsubsection*}{\subsubsection*{#2}}}{\IfNoValueTF{#2}{\subsubsection}{\subsubsection{#2}}}}
\DeclareDocumentCommand{\pr}{s o}{\IfBooleanTF{#1}{\IfNoValueTF{#2}{\paragraph*}{\paragraph*{#2}}}{\IfNoValueTF{#2}{\paragraph}{\paragraph{#2}}}}
\DeclareDocumentCommand{\spr}{s o}{\IfBooleanTF{#1}{\IfNoValueTF{#2}{\subparagraph*}{\subparagraph*{#2}}}{\IfNoValueTF{#2}{\subparagraph}{\subparagraph{#2}}}}
\DeclareDocumentCommand{\sspr}{s o}{\IfBooleanTF{#1}{\IfNoValueTF{#2}{\subsubparagraph*}{\subsubparagraph*{#2}}}{\IfNoValueTF{#2}{\subsubparagraph}{\subsubparagraph{#2}}}}
\renewcommand\theparagraph{\arabic{paragraph}.}
\renewcommand\thesubparagraph{(\arabic{subparagraph})}
\renewcommand\thesubsubparagraph{\boxed{\arabic{subsubparagraph}}}
\titlecontents{section}[2em]{\addvspace{1em}}{\thecontentslabel\quad}{}{\titlerule*[1pc]{.}\contentspage}
\titlecontents{subsection}[4em]{\addvspace{0.5em}}{\thecontentslabel\quad}{}{\titlerule*[1pc]{.}\contentspage}
\titlecontents{subsubsection}[6em]{\addvspace{0.5em}}{\thecontentslabel\ }{}{\titlerule*[1pc]{.}\contentspage}
\ifnum\value{CJK}=1\ifnum\value{SectionLanguage}=1\ifnum\value{CJKLanguage}=0
\renewcommand\thesection{第\zhnumber[style={Traditional,Normal}]{\arabic{section}}節}
\renewcommand\thesubsection{\zhnumber[style={Traditional,Normal}]{\arabic{subsection}}、}
\renewcommand\thesubsubsection{(\zhnumber[style={Traditional,Normal}]{\arabic{subsubsection}})}
\else\ifnum\value{CJKLanguage}=1
\renewcommand\thesection{第\zhnumber[style={Simplified,Normal}]{\arabic{section}}节}
\renewcommand\thesubsection{\zhnumber[style={Simplified,Normal}]{\arabic{subsection}}、}
\renewcommand\thesubsubsection{(\zhnumber[style={Simplified,Normal}]{\arabic{subsubsection}})}
\else
\renewcommand\thesection{\arabic{section}}
\renewcommand\thesubsection{\Roman{subsection}}
\renewcommand\thesubsubsection{\roman{subsubsection}}
\fi\fi
\else
\renewcommand\thesection{\arabic{section}}
\renewcommand\thesubsection{\Roman{subsection}}
\renewcommand\thesubsubsection{\roman{subsubsection}}
\fi\else
\renewcommand\thesection{\arabic{section}}
\renewcommand\thesubsection{\Roman{subsection}}
\renewcommand\thesubsubsection{\roman{subsubsection}}
\fi
\makeatletter
\@ifclassloaded{article}{}{
\titleformat{\chapter}[block]{\LARGE\bfseries}{\thechapter}{1em}{}
\DeclareDocumentCommand{\ch}{s o}{\IfBooleanTF{#1}{\IfNoValueTF{#2}{\chapter*}{\chapter*{#2}}}{\IfNoValueTF{#2}{\chapter}{\chapter{#2}}}}
\titlecontents{chapter}[0em]{\addvspace{1em}}{\thecontentslabel\quad}{}{\titlerule*[1pc]{.}\contentspage}
\ifnum\value{CJK}=1\ifnum\value{SectionLanguage}=1\ifnum\value{CJKLanguage}=0
\renewcommand\thechapter{第\zhnumber[style={Traditional,Normal}]{\arabic{chapter}}章}
\else\ifnum\value{CJKLanguage}=1
\renewcommand\thechapter{第\zhnumber[style={Simplified,Normal}]{\arabic{chapter}}章}
\else
\renewcommand\thechapter{Chapter \arabic{chapter}}
\fi\fi
\else
\renewcommand\thechapter{Chapter \arabic{chapter}}
\fi\else
\renewcommand\thechapter{Chapter \arabic{chapter}}
\fi
}
\makeatother
\ifnum\value{CJK}=1\ifnum\value{ZhRenew}=1\ifnum\value{CJKLanguage}=0
\newcommand{\temtoday}{\zhtoday}
\renewcommand{\contentsname}{\hfill\text{目錄}\hfill}
\else\ifnum\value{CJKLanguage}=1
\newcommand{\temtoday}{\zhtoday}
\renewcommand{\contentsname}{\hfill\text{目录}\hfill}
\else
\newcommand{\temtoday}{\today}
\fi\fi
\else
\newcommand{\temtoday}{\today}
\fi\else
\newcommand{\temtoday}{\today}
\fi
\newcommand{\temtitle}{\thispagestyle{empty}\Needspace{1\textheight}\maketitle\setcounter{page}{1}}
\newcommand{\temtoc}{\Needspace{1\textheight}\tableofcontents}
\newcommand{\temdoc}{\Needspace{1\textheight}\setcounter{page}{1}}
\newcommand{\titletocdoc}{\onehalfspacing\temtitle\temtoc\temdoc}
\newcommand{\titledoc}{\onehalfspacing\temtitle\temdoc}
\newcommand{\tocdoc}{\onehalfspacing\temtoc\temdoc}
\PhysicsPatch
